\section{Implications for forecasting and model evaluation}
\label{sec:discussion}

Consider a forecaster asked to predict public reaction to a policy announcement
in an unfamiliar jurisdiction. The forecaster must distinguish outcome-relevant news from
topic-matched chatter, choose the comparison population or institution that
supplies the response pattern, and revise that choice when direct local evidence
arrives. A new convention must be learned from the current case, not imported
from a familiar word. Under domain or mechanism shift, training must preserve
the relationship, not the answer distribution alone.

The prospective markets test selective updating in the wild. Coin City makes the chosen analogue observable and puts
context in conflict with local evidence. Arbitrary labels remove stable semantics;
activation patching tests one internal link; and the transfer control breaks the
episode-level correspondence without changing the answer distribution. The
historical-market study tests accuracy on real outcomes. These studies connect
selective updating, relationship choice, local revision, and transfer.

Forecasting makes conditional relationships testable because every probability
commits to how evidence, context, and outcomes connect. The relevant test is
whether that probability changes when the relationship changes.

These contrasts extend beyond forecasting. For synthetic populations, paired
relationship changes test whether conditional structure survives when marginal
averages resemble human data
\citep{park2023generative,aher2023simulate,argyle2023out,bisbee2024synthetic}.
A persuasion study holds a message fixed while changing who delivers it; a
coordination study preserves payoffs while changing what participants know
about one another; and a policy analysis preserves surface language while
changing the institutional pathway. Local
observations test correction of a misleading contextual cue, arbitrary labels
separate local induction from word association, and distribution-preserving
derangements separate relationship learning from output style or marginal
frequencies. These local behavioral tests evaluate conditional structure
directly without treating the model as a human proxy.

Parameter count and headline accuracy are insufficient model-selection
criteria. Arbitrary-cue recovery changes across populations and scales, the
controlled transfer effect replicates at two model sizes, and held-out market
performance is nonmonotonic. Validation should therefore be local to the
population, mechanism, and contrast that matter for the intended use.

\section{Limitations}
\label{sec:limitations}

The experiments test specific forms of relationship use. The real-market study
measures selective updating under natural-language packets. Coin City isolates
relationship selection between two response regimes and repeats the test in a
harder population. Future work can add exact coefficient recovery, unreliable
cues, new variables, longer horizons, and interactive multi-agent settings.

Representation analyses cover seven open checkpoints; activation patching tests
one checkpoint and one label representation. Training comparisons use three
seeds with seed-first inference. The main results use held-out tasks, paired
contrasts, and fixed estimators. Historical-market training improves every
tested base, and several checkpoints reach the crowd-level floor. A larger model
roster and later market cohorts would test the scaling pattern.

\section{Ethics, artifacts, and conclusion}
\label{sec:conclusion}

The studies use public market metadata and prices, model outputs, synthetic
worlds, and an unpaid adult review of synthetic materials. The review collected
no identifying or sensitive information and did not constitute human-subjects
research under institutional guidelines. Counterfactual news packets were used
offline and were never posted to markets or used for trading. The paper reports
aggregate model behavior, not a deployment or trading recommendation.

An anonymized artifact contains prompts, fixed manifests, saved generations,
analysis code, and figure builders.\footnote{%
\href{https://huggingface.co/datasets/anonymous-review/forecasting-study-data}%
{Study data and artifacts}; \href{https://github.com/anonymous-review/forecasting-study-code}%
{code repository}. Links anonymized for review.}
The supplementary material records censoring, estimators, missingness,
hyperparameters, and validation checks.

Across real markets and controlled tasks, LM forecasts select
context-dependent predictive relationships and revise them as direct evidence
arrives. Episode-matched training carries this relationship use across domain
and mechanism shifts at two model scales. On real markets, outcome-based
training improves every tested base, and several checkpoints reach the
crowd-level floor. These behaviors can be tested directly without assuming a
general world model.

A practical evaluation should pair outcome accuracy with targeted contrasts that
expose which relationship drives a prediction. Evaluators can hold the topic and
answer distribution fixed while changing the relevant analogue, provide local
evidence that conflicts with context, and test whether an induced relationship
survives changes in domain or mechanism. These tests expose failures hidden by an
aggregate score: reacting to lexical resemblance, preserving a population prior
when the episode differs, or carrying a valid relationship into a setting where
it no longer applies.

A useful forecasting system must meet both standards: accuracy on real outcomes
and controlled responsiveness to changes in evidence and context. The experiments
here provide a concrete way to test that requirement across the populations,
mechanisms, and decisions that define the intended use.
